\chapter{Lesson 8: Making Appointments}

\section{Vocabulary}

\begin{center}
\begin{tabular}{llll}
\toprule
Characters & Pinyin & Part of speech & Meaning \\
\midrule
\zh{喂} & wéi / wèi & interjection & hello (on the phone) \\
\zh{打电话} & dǎ diànhuà & verb-object & to make a phone call \\
\zh{给} & gěi & preposition & to, for \\
\zh{发短信} & fā duǎnxìn & verb-object & to send a text message \\
\zh{就是} & jiù shì & & it's me / precisely \\
\zh{位} & wèi & measure word & (polite, for people) \\
\zh{时间} & shíjiān & noun & time \\
\zh{问题} & wèntí & noun & question, problem \\
\zh{事(儿)} & shì(r) & noun & matter, affair \\
\zh{开会} & kāi huì & verb-object & to have a meeting \\
\zh{上课} & shàng kè & verb-object & to go to class \\
\zh{下课} & xià kè & verb-object & to get out of class \\
\zh{办公室} & bàngōngshì & noun & office \\
\zh{等} & děng & verb & to wait \\
\zh{别} & bié & adverb & don't \\
\zh{要} & yào & modal & will, be going to; need to \\
\zh{得} & děi & modal & must, have to \\
\zh{考试} & kǎo shì & verb-object / noun & to take a test; exam \\
\zh{帮} & bāng & verb & to help \\
\zh{准备} & zhǔnbèi & verb & to prepare \\
\zh{练习} & liànxí & verb / noun & to practise; exercise \\
\zh{以后} & yǐhòu & noun & after \\
\zh{以前} & yǐqián & noun & before \\
\zh{的时候} & de shíhou & & when \ldots \\
\zh{因为} & yīnwèi & conjunction & because \\
\zh{所以} & suǒyǐ & conjunction & so, therefore \\
\zh{有空(儿)} & yǒu kòng(r) & verb-object & to have free time \\
\zh{见面} & jiàn miàn & verb-object & to meet up \\
\zh{一会儿} & yíhuìr & noun & a moment \\
\bottomrule
\end{tabular}
\end{center}

\section{Phone and texting conventions}

\begin{aside}
\zh{喂} on the phone is usually said in the 2nd tone (\emph{wéi}); a sharp 4th tone
\emph{wèi} sounds like ``Hey, you!''. A caller asks \zh{请问，王老师在吗？} ``Is
Professor Wang there?''; the answer \zh{我就是。} means ``Speaking.''
By the late 2010s most students in China arranged meetings over WeChat (\zh{微信}
Wēixìn) rather than calling; voice messages (\zh{语音} yǔyīn) are as common as typed
ones.
\end{aside}

\section{Grammar}

\subsection{The preposition \zh{给}}

\zh{给} + person goes before the verb to mark the recipient of a call, text or letter.
\begin{itemize}
  \item \zh{我明天给你打电话。} I'll call you tomorrow.
  \item \zh{他常常给妈妈发短信。} He often texts his mom.
\end{itemize}
*\zh{我打电话给你} is heard in southern speech, but the textbook form puts \zh{给}
first.

\subsection{Modal verbs \zh{要} and \zh{得}}

\begin{itemize}
  \item \zh{要} ``will, be going to'' (plans): \zh{我下午要开会。} I'm going to have a
        meeting this afternoon.
  \item \zh{得} \emph{děi} ``must'' (obligation): \zh{我今天晚上得准备考试。} I have to
        prepare for an exam tonight.
\end{itemize}

\begin{important}[Negating \zh{得}]
\zh{得} \emph{děi} is not negated with \zh{不}. ``Don't have to'' is \zh{不用}
\emph{búyòng}: \zh{你不用等我。} You don't have to wait for me.
\end{important}

\subsection{Negative commands with \zh{别}}

\zh{别} + verb = ``don't''. It is softer than \zh{不要}.
\begin{itemize}
  \item \zh{别客气。} Don't be polite / make yourself at home.
  \item \zh{别等我了，我四点才下课。} Don't wait for me; I don't get out of class until four.
\end{itemize}

\subsection{Time clauses: \zh{以前}, \zh{以后}, \zh{的时候}}

Chinese time clauses end in a time word, and the whole clause comes first, before the
main clause (big-to-small again: the background before the event).

\begin{center}
\begin{tikzpicture}[x=1cm]
  \draw[thick, -{Stealth}] (0,0) -- (12,0) node[right, font=\small] {time};
  \fill[red!70!black] (6,0) circle (0.12);
  \node[above=4pt, font=\small, align=center] at (6,0.1) {X\\ \zh{下课}\\ class ends};
  \draw[thick, blue!60!black] (0.5,-0.4) -- (5.6,-0.4);
  \node[below, font=\small, align=center] at (3,-0.45) {\zh{下课以前} before class ends};
  \draw[thick, green!40!black] (6.4,-0.4) -- (11.5,-0.4);
  \node[below, font=\small, align=center] at (9,-0.45) {\zh{下课以后} after class};
  \draw[thick, orange!80!black] (5.2,-1.3) -- (6.8,-1.3);
  \node[below, font=\small] at (6,-1.35) {\zh{下课的时候} when class ends};
\end{tikzpicture}
\end{center}

\begin{itemize}
  \item \zh{下课以后我们去咖啡馆吧。} Let's go to the café after class.
  \item \zh{考试以前我得练习发音。} Before the exam I need to practise pronunciation.
  \item \zh{吃饭的时候别看手机。} Don't look at your phone while eating.
\end{itemize}

\begin{important}[Word order of time clauses]
English: ``I'll call you \emph{after I get home}.'' Chinese: \zh{我回家以后给你打电话。}
The clause with \zh{以后} comes first; it is never placed at the end.
\end{important}

\subsection{\zh{因为\ldots{}所以\ldots}}

Cause and effect are often both marked (unlike English, where ``because \ldots\ so''
is wrong). Either half may be dropped.
\begin{itemize}
  \item \zh{因为我明天有考试，所以今天晚上不能去看电影。}
  \item \zh{我今天很忙，所以不能跟你见面。}
\end{itemize}

\subsection{Asking for availability}

\begin{itemize}
  \item \zh{你今天下午有空吗？} Are you free this afternoon?
  \item \zh{你什么时候有时间？} When do you have time?
  \item \zh{我们几点见面？在哪儿见面？} What time and where shall we meet?
  \item \zh{没问题。} No problem.
  \item \zh{不好意思，我那天有事。} Sorry, I'm busy that day.
\end{itemize}

\section{Planning a week}

A sample week, used to practise the dialogue below.

\begin{center}
\begin{tabular}{lll}
\toprule
Day & Morning & Afternoon / evening \\
\midrule
\zh{星期一} & \zh{上中文课} & \zh{在图书馆看书} \\
\zh{星期二} & \zh{上数学课} & \zh{三点开会} \\
\zh{星期三} & \zh{上中文课} & \zh{有空} \\
\zh{星期四} & \zh{考试} & \zh{给妈妈打电话} \\
\zh{星期五} & \zh{有空} & \zh{跟朋友去看电影} \\
\bottomrule
\end{tabular}
\end{center}

\begin{example}[Making an appointment with the professor]\leavevmode
\begin{itemize}
  \item \textbf{A:} \zh{喂，请问，常老师在吗？}
  \item \textbf{B:} \zh{我就是。您是哪位？}
  \item \textbf{A:} \zh{老师好，我是白英爱。我想问您几个问题。您今天下午有时间吗？}
  \item \textbf{B:} \zh{今天下午我要开会。明天上午怎么样？}
  \item \textbf{A:} \zh{明天上午我有考试。您星期三下午有空吗？}
  \item \textbf{B:} \zh{星期三下午可以。两点半到我的办公室来吧。}
  \item \textbf{A:} \zh{好，谢谢老师。星期三见。}
\end{itemize}
\end{example}

\section*{Review}
\begin{itemize}
  \item \zh{给} + person + verb (\zh{给你打电话}).
  \item \zh{要} for plans; \zh{得} \emph{děi} for obligation, negated as \zh{不用}.
  \item \zh{别} + verb = don't.
  \item X \zh{以前/以后/的时候} clauses come first.
  \item \zh{因为\ldots{}所以\ldots} can be used together.
  \item Polite phone openings: \zh{喂，请问\ldots{}在吗？}---\zh{我就是}.
\end{itemize}
